\chapter{Implementazione del Framework Metodologico}
\label{ch:implementazione}

% Paragrafo introduttivo: richiama il Cap.3 (impalcatura teorica) e
% dichiara che questo capitolo ne descrive l'istanziazione concreta —
% componenti tecnici, non risultati (rimandati esplicitamente al Cap.5).
% Breve roadmap delle sezioni, sullo stesso modello dell'apertura del Cap.3.

ciao

\section{Ambiente Sperimentale e Riproducibilità}
\label{sec:ambiente_sperimentale}
% Introduzione di sezione: infrastruttura tecnica a supporto dell'intero framework.

ciao

    \subsection{Repository e Organizzazione del Codice}
    \label{subsec:repository}
    % Struttura della repository GitHub, organizzazione dei moduli/package,
    % convenzioni di versionamento, licenza, eventuale DOI/archiviazione.

    \subsection{Stack Tecnologico e Dipendenze}
    \label{subsec:stack_tecnologico}
    % Linguaggio, versione, librerie principali (es. scikit-learn, XGBoost/LightGBM,
    % shap, scipy, pandas), gestione dell'ambiente (requirements/conda), hardware
    % di esecuzione se rilevante per i tempi di training.

    \subsection{Gestione della Riproducibilità Stocastica}
    \label{subsec:riproducibilita_seed}
    % Istanziazione concreta di $s_0$ e $\mathcal{S}=\{s_1,\dots,s_K\}$
    % [rif. teorico: subsec:stabilita_stocastica] — valore numerico di $K$,
    % modalità di fissaggio dei seed nelle librerie usate, garanzie di
    % determinismo end-to-end.

\section{Caratterizzazione dei Dataset e Composizione dei Benchmark}
\label{sec:dataset_caratterizzazione}
% Introduzione: istanziazione concreta di D_S e D_T
% [rif. teorico: subsec:composizione_domini].

ciao

    \subsection{Dominio Sorgente: UNSW-NB15}
    \label{subsec:dataset_source}
    % Provenienza, dimensione, formato nativo, tassonomia nativa delle classi
    % (istanziazione concreta di $\mathcal{A}_S$), classi rimosse in
    % \textit{minority removal} e relativa soglia numerica.

    \subsection{Domini di Destinazione: Segmento Satellitare/Terrestre STIN}
    \label{subsec:dataset_target}
    % Provenienza, formato nativo, tassonomia nativa delle minacce
    % (istanziazione concreta di $\mathcal{A}_T$), classi accorpate in
    % \textit{minority merge} e criterio di raggruppamento adottato.

    \subsection{Parametrizzazione e Istanziazione dei Benchmark}
    \label{subsec:parametrizzazione_benchmark}
    % Valore numerico concreto di $\rho$, criterio without-replacement
    % applicato, conteggio effettivo delle istanze generate per ciascuna
    % delle cinque categorie [rif. teorico: subsubsec:tassonomia_configurazioni],
    % eventuale tabella riepilogativa delle cardinalità.

\section{Allineamento delle Feature e Pipeline di Preprocessing}
\label{sec:preprocessing_implementazione}
% Introduzione: istanziazione concreta dell'operatore $\psi_k$
% [rif. teorico: subsec:allineamento_feature].

ciao

    \subsection{Feature Grezze e Criteri di Filtraggio Applicati}
    \label{subsec:feature_grezze}
    % Elenco/tabella delle feature native $d_S$, $d_T$; feature scartate per
    % ciascuno dei tre criteri (identificatori, calcolabilità cross-dominio,
    % armonizzazione unità di misura); valore numerico finale di $d$.

    \subsection{Implementazione dell'Operatore di Mappatura} % $\boldsymbol{\psi}_k$
    \label{subsec:implementazione_psi}
    % Dettaglio implementativo delle trasformazioni $\psi_{k,j}$ concrete
    % (es. derivazioni, rapporti, conversioni di unità), a livello di codice
    % o pseudocodice.

    \subsection{Verifica dell'Omissione della Scalatura Esplicita}
    \label{subsec:verifica_no_scaling}
    % Nota implementativa: conferma che nessuno step di normalizzazione è
    % presente nella pipeline di training/inferenza
    % [rif. teorico: subsubsec:invarianza_monotonica], a differenza della
    % standardizzazione diagnostica (sezione successiva sugli strumenti
    % analitici).

\section{Implementazione dei Classificatori}
\label{sec:implementazione_classificatori}
% Introduzione: istanziazione concreta della tassonomia dei modelli
% [rif. teorico: sec:tassonomia_modelli].

ciao

    \subsection{Algoritmi di Base e Iperparametri}
    \label{subsec:algoritmi_iperparametri}
    % Libreria/classe concreta per ciascuna famiglia [rif. teorico:
    % subsec:famiglie_classificatori]: Random Forest, Gradient Boosting
    % (es. HistGradientBoostingClassifier/XGBoost/LightGBM); tabella degli
    % iperparametri (n. stimatori $B$, learning rate $\nu$, profondità
    % massima, criterio di split, ecc.) con eventuale strategia di tuning.

    \subsection{Istanziazione dei Modelli Aggregati}
    \label{subsec:istanziazione_aggregati}
    % Corrispondenza esplicita training-set concreto $\leftrightarrow$
    % $\mathcal{M}_{\text{base}}, \mathcal{M}_S, \mathcal{M}_H^{(T)}$
    % [rif. teorico: subsec:modelli_aggregati]; valore concreto del rapporto
    % di sparsità per $\Delta\mathcal{D}_T^{(\text{train})}$.

    \subsection{Istanziazione dei Modelli Biclasse}
    \label{subsec:istanziazione_biclasse}
    % Elenco concreto delle famiglie $a$ per cui vengono istanziati
    % $\mathcal{M}_{S,a}$ e $\mathcal{M}_{H,T,a}$
    % [rif. teorico: subsec:modelli_biclasse]; conteggio totale dei modelli
    % effettivamente addestrati.

\section{Protocollo di Addestramento e Valutazione}
\label{sec:protocollo_addestramento_valutazione}
% Introduzione: istanziazione concreta del protocollo multi-seed e della
% matrice di valutazione incrociata [rif. teorico: subsec:stabilita_stocastica].

ciao

    \subsection{Orchestrazione del Ciclo di Addestramento Multi-Seed}
    \label{subsec:orchestrazione_training}
    % Flusso implementativo concreto: come vengono eseguiti $K$ run per
    % ciascun modello, gestione della persistenza dei modelli congelati,
    % eventuale parallelizzazione/tempi di calcolo.

    \subsection{Calcolo delle Metriche e Costruzione della Matrice di Valutazione} % $H_k$
    \label{subsec:calcolo_metriche_hk}
    % Implementazione concreta di TPR/TNR/Precision/F1
    % [rif. teorico: subsubsec:metriche_aggregazione], libreria usata
    % (es. scikit-learn.metrics), formato di persistenza della matrice
    % (CSV/DataFrame), eventuale esportazione delle heatmap.

    \subsection{Implementazione del Test di Welch}
    \label{subsec:implementazione_welch}
    % Libreria concreta (es. scipy.stats.ttest\_ind con equal\_var=False)
    % [rif. teorico: subsubsec:welch\_test], soglia $\alpha$ adottata,
    % modalità di reporting dei $p$-value.

\section{Strumenti di Analisi Diagnostica e Spiegabilità}
\label{sec:strumenti_analisi_diagnostica}
% Introduzione: istanziazione concreta del framework analitico
% [rif. teorico: sec:framework_analitico].

ciao

    \subsection{Standardizzazione Diagnostica, PCA e KDE}
    \label{subsec:implementazione_pca_kde}
    % Libreria concreta (es. sklearn.decomposition.PCA,
    % scipy.stats.gaussian\_kde o seaborn), parametro di bandwidth
    % effettivamente usato/selezionato per $h_{\text{KDE}}$
    % [rif. teorico: subsubsec:kde\_features], formato dei grafici prodotti.

    \subsection{Implementazione di SHAP e TreeSHAP}
    \label{subsec:implementazione_shap}
    % Libreria concreta (\texttt{shap}), classe \texttt{TreeExplainer},
    % dimensione concreta del campione stratificato
    % $|\mathcal{D}_{\text{sample}}^{(\text{test})}|$
    % [rif. teorico: subsubsec:treeshap], tempi/costo computazionale.

    \subsection{Pipeline di Aggregazione SHAP--KDE}
    \label{subsec:pipeline_shap_kde}
    % Implementazione concreta del contatore $C_j$ e della selezione
    % delle tre feature finali [rif. teorico: subsubsec:aggregazione_shap_kde].

\section{Sintesi della Pipeline Implementativa}
\label{sec:sintesi_pipeline_implementativa}
% Chiusura del capitolo: eventuale diagramma/tabella riepilogativa
% end-to-end (dataset -> preprocessing -> training multi-seed -> analisi),
% che funge da raccordo esplicito verso il Capitolo 5 (risultati).

ciao